\subsection{Rekomendacje i wyniki MCDA}
Repozytorium nie zawiera zapisanego logu pełnej symulacji z liczbowymi
wynikami poszczególnych ofert. Dlatego poniższe szczegółowe dane
pochodzą z deterministycznych przypadków testowych, a nie z przebiegu
80 samochodów. Wyniki całej symulacji podano osobno w tabeli~\ref{tab:runs}.

W \code{RecommendationServiceTests} cztery zgodne oferty SUV kosztują
łącznie 120\,000, 130\,000, 140\,000 i 150\,000 PLN.
Każda ma koszt dodatkowy 2\,000 PLN, rok 2024, pojemność 1998 cm$^3$,
paliwo PETROL i spalanie 7,0 l/100 km. Oferta SEDAN za 90\,000 PLN
jest odrzucana mimo niższej ceny.

\begin{table}[htbp]
\centering
\begin{tabular}{lrrl}
\toprule
Oferta & Koszt całkowity [PLN] & Wynik MCDA & Decyzja\\
\midrule
CAR-001 & 120\,000 & 1,00 & TOP 1\\
CAR-002 & 130\,000 & 0,85 & TOP 2\\
CAR-003 & 140\,000 & 0,70 & TOP 3\\
CAR-004 & 150\,000 & 0,55 & Poza TOP 3\\
CAR-005 & 90\,000 & --- & Niezgodne nadwozie\\
\bottomrule
\end{tabular}
\caption{Ranking testowego zbioru ofert. Oceny liczbowe wyliczono
z danych testu i wzoru implementacji; asercje testu sprawdzają
kolejność TOP~3 oraz zakres ocen, nie te konkretne wartości.}
\end{table}

Cena i koszt całkowity odpowiadają tu łącznie za 45\% oceny.
Pozostałe cechy są jednakowe, więc każda oferta uzyskuje za nie 0,55.
Dla CAR-002 normalizacja ceny oraz kosztu daje $2/3$, stąd
$S=0{,}55+0{,}45\cdot2/3=0{,}85$.
Test remisu przy jednakowych parametrach preferuje Seller02
z reputacją 5,0 przed Seller01 z reputacją 3,0.

\subsection{Negocjacje}
W \code{NegotiationServiceTests} samochód ma cenę 145\,000 PLN
i koszt dodatkowy 3\,000 PLN. Minimalny koszt po rabacie wynosi
$145\,000\cdot0{,}85+3\,000=126\,250$ PLN.
Oferta kupującego 119\,000 PLN przy limicie 148\,000 PLN prowadzi
do kontrpropozycji i akceptacji 133\,500 PLN już w pierwszej rundzie.
Oszczędność względem katalogowego kosztu 148\,000 PLN wynosi
14\,500 PLN. Drugi przypadek, z limitem 120\,000 PLN,
kończy się odrzuceniem po trzech rundach.

\subsection{Reputacja i trendy}
Test reputacji zapisuje dla Seller01 oceny 5 i 4 z dwóch transakcji.
Uzyskana średnia to $(5+4)/2=4{,}5$, a licznik ocen wynosi 2.
Ponowna ocena tej samej transakcji jest odrzucana.

Test trendów przetwarza cztery transakcje o wartościach 100\,000,
105\,000, 130\,000 i 140\,000 PLN. Wynik obejmuje cztery sprzedaże,
średnią 118\,750 PLN, najpopularniejsze nadwozie SUV, model Toyota RAV4,
paliwo PETROL i kierunek RISING. Są to agregaty kontrolowanego
zbioru testowego, a nie estymacja trendu rzeczywistego rynku.

\subsection{Elasticsearch i Kibana}
Plik \code{kibana/dashboard.ndjson} zawiera dwa widoki danych, pięć
wizualizacji Vega i dashboard \emph{SSISW --- giełda samochodowa}.
Wizualizacje przedstawiają liczbę pojazdów według marki, nadwozia
i paliwa, rozkład roku produkcji oraz transakcje według statusu
COMPLETED / RETURNED. Odwołują się do indeksów \code{cars}
i \code{transactions}; filtr \code{runId} pozwala wybrać przebieg.
Zrzuty działającego dashboardu nie zostały dostarczone z materiałami,
więc ich nie zastępowano przykładowymi obrazami.
